Écris une fonction \texttt{tri\_selection(tableau)} qui trie \emph{en place} (le tableau est modifié) une liste d'entiers \texttt{tableau}, par ordre croissant.

On utilise l'algorithme suivant : on parcourt le tableau de gauche à droite ; à chaque position, on cherche le minimum du tableau entre cette position et la fin, puis on échange ces deux valeurs.

Exemples : après \texttt{tri\_selection(tab)} avec \texttt{tab = [1, 52, 6, -9, 12]}, \texttt{tab} vaut \texttt{[-9, 1, 6, 12, 52]} ; une liste vide reste vide ; \texttt{[9]} reste \texttt{[9]}.
